%-------------------------
% Resume in Latex
% Author : Jake Gutierrez
% Based off of: https://github.com/sb2nov/resume
% License : MIT
%------------------------

\documentclass[letterpaper,11pt]{article}

\usepackage{latexsym}
\usepackage[empty]{fullpage}
\usepackage{titlesec}
\usepackage{marvosym}
\usepackage[usenames,dvipsnames]{color}
\usepackage{verbatim}
\usepackage{enumitem}
\usepackage[hidelinks]{hyperref}
\usepackage{fancyhdr}
\usepackage[english]{babel}
\usepackage{tabularx}
\input{glyphtounicode}


%----------FONT OPTIONS----------
% sans-serif
% \usepackage[sfdefault]{FiraSans}
% \usepackage[sfdefault]{roboto}
% \usepackage[sfdefault]{noto-sans}
% \usepackage[default]{sourcesanspro}

% serif
% \usepackage{CormorantGaramond}
% \usepackage{charter}


\pagestyle{fancy}
\fancyhf{} % clear all header and footer fields
\fancyfoot{}
\renewcommand{\headrulewidth}{0pt}
\renewcommand{\footrulewidth}{0pt}

% Adjust margins
\addtolength{\oddsidemargin}{-0.5in}
\addtolength{\evensidemargin}{-0.5in}
\addtolength{\textwidth}{1in}
\addtolength{\topmargin}{-.5in}
\addtolength{\textheight}{1.0in}

\urlstyle{same}

\raggedbottom
\raggedright
\setlength{\tabcolsep}{0in}

% Sections formatting
\titleformat{\section}{
  \vspace{-4pt}\scshape\raggedright\large
}{}{0em}{}[\color{black}\titlerule \vspace{-5pt}]

% Ensure that generate pdf is machine readable/ATS parsable
\pdfgentounicode=1

%-------------------------
% Custom commands
\newcommand{\resumeItem}[1]{
  \item\small{
    {#1 \vspace{-2pt}}
  }
}

\newcommand{\resumeSubheading}[4]{
  \vspace{-2pt}\item
    \begin{tabular*}{0.97\textwidth}[t]{l@{\extracolsep{\fill}}r}
      \textbf{#1} & #2 \\
      \textit{\small#3} & \textit{\small #4} \\
    \end{tabular*}\vspace{-7pt}
}

\newcommand{\resumeSubSubheading}[2]{
    \item
    \begin{tabular*}{0.97\textwidth}{l@{\extracolsep{\fill}}r}
      \textit{\small#1} & \textit{\small #2} \\
    \end{tabular*}\vspace{-7pt}
}

\newcommand{\resumeProjectHeading}[2]{
    \item
    \begin{tabular*}{0.97\textwidth}{l@{\extracolsep{\fill}}r}
      \small#1 & #2 \\
    \end{tabular*}\vspace{-7pt}
}

\newcommand{\resumeSubItem}[1]{\resumeItem{#1}\vspace{-4pt}}

\renewcommand\labelitemii{$\vcenter{\hbox{\tiny$\bullet$}}$}

\newcommand{\resumeSubHeadingListStart}{\begin{itemize}[leftmargin=0.15in, label={}]}
\newcommand{\resumeSubHeadingListEnd}{\end{itemize}}
\newcommand{\resumeItemListStart}{\begin{itemize}}
\newcommand{\resumeItemListEnd}{\end{itemize}\vspace{-5pt}}

%-------------------------------------------
%%%%%%  RESUME STARTS HERE  %%%%%%%%%%%%%%%%%%%%%%%%%%%%


\begin{document}

%----------HEADING----------
% \begin{tabular*}{\textwidth}{l@{\extracolsep{\fill}}r}
%   \textbf{\href{http://sourabhbajaj.com/}{\Large Sourabh Bajaj}} & Email : \href{mailto:sourabh@sourabhbajaj.com}{sourabh@sourabhbajaj.com}\\
%   \href{http://sourabhbajaj.com/}{http://www.sourabhbajaj.com} & Mobile : +1-123-456-7890 \\
% \end{tabular*}

\begin{center}
    {\Huge \bfseries \rmfamily Shams Mohammed Saiham Kabir} \\[1pt]  % Serif font (Times), bold
    \vspace*{0.35cm}  % Adjust the value as needed
    \small 
    \bfseries +49 15751104740 \hspace{4pt} | \hspace{4pt} 
    \href{mailto:shams@example.com}{\bfseries shamsmohammed726@gmail.com} \hspace{4pt} | \hspace{4pt}
    \href{https://shams72.github.io/reactPortfolio/}{\bfseries shams-msk.com}
\end{center}



%-----------EDUCATION-----------
\section{\textbf{\fontsize{13.5}{22}\rmfamily Ausbildung}}  
  \resumeSubHeadingListStart
     \resumeSubheading
 {\textbf{Hochschule Darmstadt}}{\textbf{Okt. 2022 -- heute}}
  {\textbf{Bachelor of Science in Informatik}}{}
      \resumeItemListStart
        \resumeItem{Relevante Kurse: Software Engineering, Webentwicklung, Künstliche Intelligenz, Rechnernetze, Betriebssysteme, Algorithmen und Datenstrukturen.}
        \resumeItem{Voraussichtlicher Abschluss: März 2026.}
      \resumeItemListEnd
  \resumeSubHeadingListEnd


%-----------EXPERIENCE-----------

\section{\textbf{\fontsize{13.5}{18}\rmfamily Berufserfahrung}} 
  \resumeSubHeadingListStart

    \resumeSubheading
      {abasoft EDV-Programme GmbH}{\textbf{Nov. 2024 -- Heute}}
      {\textbf{Werkstudent im Bereich Softwareentwicklung}}{}
      \vspace*{0.05cm}
      \resumeItemListStart
        \resumeItem{Entwicklung und Implementierung neuer Funktionen basierend auf den Anforderungen mit PHP, HTML, CSS und JavaScript}
        \resumeItem{Fehlerbehebung und Erweiterung von Funktionen in enger Zusammenarbeit mit dem Entwicklungsteam mittels React}
      \resumeItemListEnd

    \resumeSubheading
      {Hochschule Darmstadt}{\textbf{Apr. 2024 -- Heute}}
      {\textbf{{Studentische Hilfskraft für ein Forschungsprojekt}}}{}
      \vspace*{0.05cm}
      \resumeItemListStart
        \resumeItem{Erstellung von Mock-Dashboards mit Python unter Verwendung der Bibliotheken Streamlit, Voila, etc.}
        \resumeItem{Containerisierung von Python mit Docker}
      \resumeItemListEnd

  \resumeSubHeadingListEnd

\section{\textbf{\fontsize{13.5}{18}\rmfamily Technische Fähigkeiten}} 
 \begin{itemize}[leftmargin=0.15in, label={}]
     \resumeSubHeadingListStart
      \resumeItem{\textbf{Sprachen:}}{PHP, JavaScript, TypeScript, HTML/CSS, Python, C++}
      \resumeItem{\textbf{Frameworks:}}{React, Qt GUI}
      \resumeItem{\textbf{Datenbankmanagement:}}{SQL (PostgreSQL), db4o, MongoDB}
      \resumeItem{\textbf{Entwicklungswerkzeuge:}}{Git, Docker}
    \resumeSubHeadingListEnd
\end{itemize}
      
% -----------Multiple Positions Heading-----------
%    \resumeSubSubheading
%     {Software Engineer I}{Oct 2014 - Sep 2016}
%     \resumeItemListStart
%        \resumeItem{Apache Beam}
%          {Apache Beam is a unified model for defining both batch and streaming data-parallel processing pipelines}
%     \resumeItemListEnd
%    \resumeSubHeadingListEnd
%-------------------------------------------


%-----------PROJECTS-----------  
\section{\textbf{\fontsize{13.5}{18}\rmfamily Projekte}} 
    \resumeSubHeadingListStart
    
      \resumeProjectHeading
          {\textbf{ShopEsy} $|$ \textbf{React, TypeScript, Express, Jest, Supertest}}{\href{https://github.com/shams72/ShopEsy} {\textbf{[GitHub Link]}}}
          \resumeItemListStart
                   \resumeItem{Entwicklung einer vollständigen CRUD-Anwendung zur Verwaltung von Einkaufsartikeln in Listen, einschließlich Funktionen zum Erstellen, Anzeigen, Bearbeiten und Löschen von Einträgen.}

          \resumeItemListEnd

      \resumeProjectHeading
          {\textbf{Yapple Social Media} $|$ \textbf{React, TypeScript, Express, MongoDB, Jest, Supertest}}{\href{https://github.com/shams72/Yapple-Social} {\textbf{[GitHub Link]}}}
          \resumeItemListStart
              \resumeItem{Teamorientierte Entwicklung einer Social-Media-Plattform im Rahmen eines Webentwicklungsprojekts, mit Funktionen zum Erstellen von Beiträgen, Communities und zur Interaktion zwischen Nutzern.}
        \resumeItem{Verantwortlich für die Implementierung neuer Features im Frontend und Backend gemäß projektbezogenen Anforderungen.}
          \resumeItemListEnd

      \resumeProjectHeading
          {\textbf{Reisebüro-Anwendung} $|$ {\textbf{Qt Gui, C++}}}{\href{https://github.com/shams72/TravelAgency1} {\textbf{[GitHub Link]}}}
          \resumeItemListStart
            \resumeItem{Entwicklung eines Reisebürosystems mit Qt-Bibliotheken.}
                   \resumeItem{Eigenständige Vertiefung in die C++- und Qt-Entwicklung durch Nutzung relevanter Online-Ressourcen und Tutorials.}

          \resumeItemListEnd

      \resumeProjectHeading
          {\textbf{FunTaskTic To-Do-App} $|$ \textbf{Java, XML, Android Studio}}{\href{https://github.com/shams72/FunTaskTic} {\textbf{[GitHub Link]}}}
          \resumeItemListStart
            \resumeItem{Gemeinsame Entwicklung einer To-Do-App namens „FunTaskTic“ in Java unter Verwendung von Android Studio.}
            \resumeItem{Integration von Funktionen basierend auf umfassender Benutzerforschung durch das Team, um eine benutzerzentrierte Anwendung zu gewährleisten.}
          \resumeItemListEnd

      \resumeProjectHeading
          {\textbf{Skyrim Website} $|$ \textbf{React, TypeScript, MongoDB, Axios, Express, Supertest}}{\href{https://github.com/shams72/Skyrim} {\textbf{[GitHub Link]}}}
         \resumeItemListStart
            {\fontsize{12}{11}\selectfont
            \resumeItem{Gemeinsame Entwicklung einer Kartenanwendung für Skyrim mit einem Navigationsdienst, entwickelt mit DevOps-Prinzipien im Rahmen des Moduls „Projekt System Entwicklung“.}
        \resumeItem{Implementierung und Test von API-Endpunkten unter Verwendung von Jest und Supertest.}
            }
        \resumeItemListEnd
        

\resumeSubHeadingListEnd
%-----------PROGRAMMING SKILLS-----------



%



%-------------------------------------------
\end{document}
